\chapter{Mười Lát Cắt Lấy Lại Đời Sống}
\markboth{Mười Lát Cắt Lấy Lại Đời Sống}{Mười Lát Cắt Lấy Lại Đời Sống}

\latcat{1. Một Bữa Cơm Sạch Màn Hình}{Chỉ một bữa cơm không điện thoại cũng đủ làm nhiều gia đình lúng túng vì quên mất cách gặp nhau.\\
\textit{(Just one screen-free meal can make many families awkward, because they have forgotten how to meet one another.)}\\
Chính sự lúng túng đó cho thấy việc làm này đáng cứu đến mức nào.\\
\textit{(That awkwardness itself proves how worth saving this practice is.)}}{Discomfort at first is often a sign of healing.}

\latcat{2. Mười Phút Chờ Mà Không Lướt}{Thử chờ mười phút mà không mở máy.\\
\textit{(Try waiting ten minutes without opening the phone.)}\\
Nếu thấy bứt rứt, đó không phải vì đời quá rỗng. Có thể là vì mình đã quá quen sống như một người bị giật dây.\\
\textit{(If restlessness comes, it may not be because life is empty. It may be because we have become too used to living like something pulled by strings.)}}{Restlessness reveals dependence.}

\latcat{3. Cất Máy Khi Người Khác Đang Nói}{Một động tác úp màn hình xuống bàn đôi khi có giá trị hơn rất nhiều câu xin lỗi vì phân tâm.\\
\textit{(Turning the screen face down can mean more than many apologies for distraction.)}\\
Vì nó nói rằng con người trước mặt mình đang được ưu tiên thật.\\
\textit{(It says the person in front of us is truly being prioritized.)}}{Respect is often visible in tiny gestures.}

\latcat{4. Tách Học Và Lướt}{Một giờ học sâu mà máy nằm xa đôi khi cứu được đầu óc hơn cả ba giờ ngồi cạnh nó với cảm giác mình vẫn đang học.\\
\textit{(One deep hour of study with the phone far away can save the mind more than three hours spent beside it while pretending to study.)}\\
Học cần kỷ luật cơ học trước khi nói đến lý tưởng tinh thần.\\
\textit{(Learning needs mechanical discipline before it can speak of spiritual ideals.)}}{Distance from distraction is often the first form of wisdom.}

\latcat{5. Đi Bộ Để Đi Bộ}{Đi bộ mà không nghe gì, không xem gì, không trả lời gì ban đầu sẽ rất lạ.\\
\textit{(Walking without listening, watching, or replying feels strange at first.)}\\
Nhưng chính cái lạ đó có thể mở lại cho đầu óc một cánh cửa đã bị niêm phong quá lâu.\\
\textit{(But that strangeness can reopen a door in the mind that has been sealed for too long.)}}{Unmediated reality feels unfamiliar only because we abandoned it.}

\latcat{6. Không Chạy Xe Với Màn Hình}{Một người tự trọng không lấy mạng mình và mạng người khác ra đánh cược với một thông báo.\\
\textit{(A person with self-respect does not gamble human life against a notification.)}\\
Cất điện thoại khi lái xe không chỉ là kỹ năng. Nó là phẩm chất.\\
\textit{(Putting away the phone while driving is not only skill. It is character.)}}{Safety is a moral decision before it is a technical one.}

\latcat{7. Trả Lại Phòng Ngủ Cho Giấc Ngủ}{Mang điện thoại lên giường là mang cả đám đông, quảng cáo, tranh cãi, cám dỗ và tiếng động vào chỗ đáng lẽ phải yên nhất.\\
\textit{(Bringing a phone to bed means bringing the crowd, ads, quarrels, temptations, and noise into the place that should be quietest.)}\\
Muốn ngủ sâu hơn, nhiều khi phải bắt đầu từ sự tàn nhẫn nhỏ: đừng cho màn hình lên giường nữa.\\
\textit{(To sleep more deeply, one often begins with a small cruelty: stop letting the screen onto the bed.)}}{Mercy to sleep often requires firmness toward the phone.}

\latcat{8. Dành Lại Chút Chán}{Người biết chịu chán một chút thường đọc được sâu hơn, nghĩ được lâu hơn, và đỡ bị nội dung rẻ tiền dắt đi hơn.\\
\textit{(People who can endure a little boredom often read more deeply, think longer, and get dragged around less by cheap content.)}\\
Chán không phải lúc nào cũng là kẻ thù. Nó có thể là cửa vào chiều sâu.\\
\textit{(Boredom is not always the enemy. It can be an entrance to depth.)}}{Boredom can be compost for thought.}

\latcat{9. Quy Tắc Ít Nhưng Giữ Thật}{Không cần lập hai mươi quy tắc rồi phá hết mười chín.\\
\textit{(There is no need to create twenty rules and break nineteen.)}\\
Chỉ cần vài luật nhỏ mà giữ thật: không máy trên bàn ăn, không máy khi chạy xe, không máy trong giờ học sâu, không máy lúc người khác đang cần mình có mặt.\\
\textit{(A few small rules kept seriously are enough: no phone at the table, none while driving, none during deep study, none when someone needs your presence.)}}{A few kept rules beat a long list of decorative intentions.}

\latcat{10. Sống Lại Không Có Nghĩa Là Về Thời Không Có Công Nghệ}{Không ai bắt con người chối bỏ tiến bộ.\\
\textit{(No one asks human beings to reject progress.)}\\
Điều cần là đặt công nghệ về đúng ngôi của nó: làm đầy tay thì được, nhưng đừng cho nó nuốt mất mắt, mất đầu, mất tim và mất đời sống.\\
\textit{(What is needed is to return technology to its proper place: let it fill the hand if necessary, but do not let it swallow the eyes, the mind, the heart, and life itself.)}}{Tools belong in the hand, not on the throne.}

\ngancach